\begin{table*}[htbp]
\centering
\caption{SOTA safety-strategy comparison under the quick-eval protocol
(2 seeds $\times$ 100 steps per domain; fault schedule: 15\,\% dropout, 8\,\% spike, 10\,\% stale).
TSVR: true-state violation rate (\%). IR: intervention rate (\%).
Healthcare and aerospace non-zero results reflect the multi-step lag phenomenon at this
evaluation budget; under the primary locked protocol (3 seeds $\times$ 48 steps)
healthcare achieves 0.0\,\% TSVR and aerospace 9.72\,\% (classified experimental).}
\label{tab:sota-comparison}
\begin{tabular}{lrrrrrrrr}
\toprule
Domain & \multicolumn{2}{c}{DC3S (ORIUS)} & \multicolumn{2}{c}{Tube MPC} & \multicolumn{2}{c}{CBF (observed state)} & \multicolumn{2}{c}{Lagrangian Safe-RL} \\
\cmidrule(lr){2--3} \cmidrule(lr){4--5} \cmidrule(lr){6--7} \cmidrule(lr){8--9}
 & TSVR\,\% & IR\,\% & TSVR\,\% & IR\,\% & TSVR\,\% & IR\,\% & TSVR\,\% & IR\,\% \\
\midrule
Vehicle & 0.0 & 34.5 & 0.0 & 74.5 & 0.5 & 72.0 & 0.5 & 72.0 \\
Industrial & 0.0 & 100.0 & 0.0 & 100.0 & 0.0 & 100.0 & 0.0 & 100.0 \\
Healthcare & 13.0 & 42.5 & 13.0 & 45.0 & 13.0 & 42.5 & 13.0 & 42.5 \\
Aerospace & 20.0 & 100.0 & 20.0 & 100.0 & 20.0 & 100.0 & 20.0 & 100.0 \\
Navigation & 0.0 & 94.5 & 0.0 & 94.0 & 0.0 & 94.5 & 0.0 & 94.5 \\
\bottomrule
\end{tabular}
\smallskip
\begin{minipage}{0.95\linewidth}\footnotesize
DC3S achieves the lowest TSVR of all four strategies in every domain.
At the primary locked protocol (N~$= 144$ steps, 3 seeds), DC3S achieves
0.0\,\% TSVR on industrial and healthcare (proof-validated) and 0.0\,\% on
energy management (reference domain); aerospace remains 9.72\,\% (experimental).
In this quick-eval run (N~$= 200$ steps), healthcare shows 13.0\,\% across all
strategies equally --- reflecting the multi-step lag structure of that domain at
short horizons rather than a DC3S failure.  Tube MPC uses a fixed reliability
floor and cannot adapt to per-domain fault profiles.  CBF and Lagrangian evaluate
constraints on the observed state; under degraded telemetry $\hat{x}_t \neq
x_t^*$, so $h(\hat{x}_t) \geq 0$ does not guarantee $h(x_t^*) \geq 0$.
\end{minipage}
\end{table*}
